% TODO make hw stylesheet
\documentclass[11pt]{scrartcl}
\usepackage[a4paper, total={6in, 8in}]{geometry}
\usepackage[usenames,dvipsnames,svgnames]{xcolor}
\usepackage[shortlabels]{enumitem}
\usepackage[framemethod=TikZ]{mdframed}
\usepackage{amsmath,amssymb,amsthm}
\usepackage[colorlinks]{hyperref}
\usepackage{multicol}
\usepackage{tabularx}
\usepackage{stmaryrd}

\hypersetup{
	colorlinks=true,
	linkcolor=blue,
	filecolor=magenta,
	urlcolor=magenta,
	pdftitle={MAT 548 Notes}
}

\usepackage{microtype}
\usepackage{mathtools}
\usepackage[headsepline]{scrlayer-scrpage}
\usepackage{thmtools}
\usepackage[textsize=footnotesize]{todonotes}

\mdfdefinestyle{mdbluebox}{roundcorner=10pt,innerbottommargin=9pt,
	linecolor=blue,backgroundcolor=TealBlue!5,}
\declaretheoremstyle[headfont=\sffamily\bfseries\color{MidnightBlue},
mdframed={style=mdbluebox},]{thmbluebox}

\mdfdefinestyle{mdbloobox}{frametitlefont=\bfseries,innerbottommargin=8pt,
	nobreak=true,backgroundcolor=TealBlue!5,linecolor=blue,}
\declaretheoremstyle[headfont=\bfseries\color{MidnightBlue},
mdframed={style=mdbloobox},headpunct={\\[3pt]},postheadspace=0pt,]{thmbloobox}

\mdfdefinestyle{mdredbox}{frametitlefont=\bfseries,innerbottommargin=8pt,
	nobreak=true,backgroundcolor=Salmon!5,linecolor=RawSienna,}
\declaretheoremstyle[headfont=\bfseries\color{RawSienna},
mdframed={style=mdredbox},headpunct={\\[3pt]},postheadspace=0pt,]{thmredbox}

\mdfdefinestyle{mdgreenbox}{linecolor=ForestGreen,backgroundcolor=ForestGreen!5,
	linewidth=2pt,rightline=false,leftline=true,topline=false,bottomline=false,}
\declaretheoremstyle[headfont=\bfseries\sffamily\color{ForestGreen!70!black},
mdframed={style=mdgreenbox},headpunct={ --- },]{thmgreenbox}

\mdfdefinestyle{mdorangebox}{linecolor=RedOrange,backgroundcolor=RedOrange!5,
	linewidth=2pt,rightline=false,leftline=true,topline=false,bottomline=false,}
\declaretheoremstyle[headfont=\bfseries\sffamily\color{RedOrange!70!black},
mdframed={style=mdorangebox},headpunct={ --- },]{thmorangebox}

\mdfdefinestyle{mdblackbox}{linecolor=black,backgroundcolor=RedViolet!5!gray!5,
	linewidth=3pt,rightline=false,leftline=true,topline=false,bottomline=false,}
\declaretheoremstyle[mdframed={style=mdblackbox}]{thmblackbox}

\declaretheoremstyle[spaceabove=6pt,spacebelow=6pt]{basehead}

\declaretheorem[style=thmbluebox,name=Theorem,numberwithin=section]{theorem}
\declaretheorem[style=thmredbox,name=Example,sibling=theorem]{example}
\declaretheorem[style=thmbluebox,name=Corollary,sibling=theorem]{corollary}
\declaretheorem[style=thmbluebox,name=Proposition,sibling=theorem]{prop}
\declaretheorem[style=thmbluebox,name=Lemma,numberwithin=theorem]{lemma}
\declaretheorem[style=thmbluebox,name=Theorem,numbered=no]{theorem*}
\declaretheorem[style=thmbluebox,name=Lemma,numbered=no]{lemma*}
\declaretheorem[style=thmgreenbox,name=Claim,numberwithin=theorem]{claim}
\declaretheorem[style=thmgreenbox,name=Claim,numbered=no]{claim*}
\declaretheorem[style=thmblackbox,name=Remark,numberwithin=theorem]{remark}
\declaretheorem[style=thmblackbox,name=Remark,numbered=no]{remark*}
\declaretheorem[style=thmgreenbox,name=Definition,sibling=theorem]{definition}
\declaretheorem[style=thmgreenbox,name=Definition,numbered=no]{definition*}
\declaretheorem[style=thmorangebox,name=Warning,sibling=theorem]{warning}
\declaretheorem[style=thmorangebox,name=Warning,numbered=no]{warning*}
\declaretheorem[style=thmblackbox,name=Exercise,sibling=theorem]{exercise}
\declaretheorem[style=thmblackbox,name=Exercise,numbered=no]{exercise*}

\addtolength{\textheight}{3.14cm}
\providecommand{\ol}{\overline}
\providecommand{\ceil}[1]{\left\lceil #1 \right\rceil}
\providecommand{\floor}[1]{\left\lfloor #1 \right\rfloor}
\newcommand{\dd}{\mathrm{d}}
\providecommand{\ul}{\underline}
\providecommand{\eps}{\varepsilon}
\providecommand{\half}{\frac{1}{2}}
\providecommand{\dang}{\measuredangle} %% Directed angle
\providecommand{\CC}{\mathbb C}
\providecommand{\PP}{\mathbb P}
\providecommand{\FF}{\mathbb F}
\providecommand{\NN}{\mathbb N}
\providecommand{\QQ}{\mathbb Q}
\providecommand{\RR}{\mathbb R}
\providecommand{\ZZ}{\mathbb Z}
\providecommand{\dg}{^\circ}
\providecommand{\ii}{\item}
\providecommand{\setdiff}{\smallsetminus}
\providecommand{\alert}[1]{{\sffamily\textbf{\textcolor{blue}{#1}}}}
\providecommand{\inv}{^{-1}}
\providecommand{\mb}[1]{\mathbf{#1}}


\reversemarginpar

\newenvironment{soln}{\begin{proof}[Solution]}{\end{proof}}
\newenvironment{walkthrough}{\noindent \textbf{Walkthrough.}}{\medskip}
\newlist{walk}{enumerate}{3}
\setlist[walk]{label=\bfseries (\alph*)}

\providecommand{\pow}{\operatorname{Pow}}
\providecommand{\Gr}{\operatorname{Gr}}
\providecommand{\vocab}[1]{{\textbf{\textcolor{ForestGreen}{#1}}}}
\providecommand{\lcm}{\operatorname{lcm}}
\providecommand{\abs}[1]{\left|#1\right|}
\providecommand{\norm}[1]{\left\|#1\right\|}
\providecommand{\floor}[1]{\left\lfloor #1\right\rfloor}
\providecommand{\ceil}[1]{\left\lceil #1\right\rceil}

\usepackage{asymptote}
\begin{asydef}
	defaultpen(fontsize(10pt));
	unitsize(1cm);
	size(8cm); // set a reasonable default
	usepackage("amsmath");
	usepackage("amssymb");
	settings.tex="pdflatex";
	settings.outformat="pdf";
	// Replacement for olympiad+cse5 which is not standard
	import geometry;
	// recalibrate fill and filldraw for conics
	void filldraw(picture pic = currentpicture, conic g, pen fillpen=defaultpen, pen drawpen=defaultpen)
	{ filldraw(pic, (path) g, fillpen, drawpen); }
	void fill(picture pic = currentpicture, conic g, pen p=defaultpen)
	{ filldraw(pic, (path) g, p); }
	// some geometry
	pair foot(pair P, pair A, pair B) { return foot(triangle(A,B,P).VC); }
	pair orthocenter(pair A, pair B, pair C) { return orthocentercenter(A,B,C); }
	pair centroid(pair A, pair B, pair C) { return (A+B+C)/3; }
	// cse5 abbreviations
	path CP(pair P, pair A) { return circle(P, abs(A-P)); }
	path CR(pair P, real r) { return circle(P, r); }
	pair IP(path p, path q) { return intersectionpoints(p,q)[0]; }
	pair OP(path p, path q) { return intersectionpoints(p,q)[1]; }
	path Line(pair A, pair B, real a=0.6, real b=a) { return (a*(A-B)+A)--(b*(B-A)+B); }
	// cse5 more useful functions
	picture CC() {
		picture p=rotate(0)*currentpicture;
		currentpicture.erase();
		return p;
	}
	pair MP(Label s, pair A, pair B = plain.S, pen p = defaultpen) {
		Label L = s;
		L.s = "$"+s.s+"$";
		label(L, A, B, p);
		return A;
	}
	pair Drawing(Label s = "", pair A, pair B = plain.S, pen p = defaultpen) {
		dot(MP(s, A, B, p), p);
		return A;
	}
	path Drawing(path g, pen p = defaultpen, arrowbar ar = None) {
		draw(g, p, ar);
		return g;
	}
\end{asydef}

\usepackage{mathrsfs}

\ihead{\footnotesize MAT 548 Notes}
\ohead{\footnotesize \today}

\title{\vspace{-2em}MAT 548: Geometry and Topology I}
\author{\large Aditya Pahuja}
\date{}
\begin{document}
\maketitle
\tableofcontents

\pagebreak

\section{August 26}\label{sec:aug26}

\begin{definition}[Topological manifold]
    \label{def:topman}
    An \vocab{$n$-dimensional topological manifold}
    is a second-countable Hausdorff space $X$ that is locally Euclidean,
    i.e. for all $x\in X$, there is an open $U\subseteq X$ containing $x$
    that is homeomorphic to an open subset of $\RR^n$.
\end{definition}

Note that this implies that manifolds are locally compact
because open sets in $\RR^n$ are locally compact.
Topological manifolds also satisfy a useful property called paracompactness:
\begin{definition}[Paracompactness]
    \label{def:paracompact}
    A topological space $X$ is paracompact if
    any open cover $\{U_\alpha\}_{\alpha\in A}$ of $X$
    has a locally finite refinement;
    this means that there is an open cover $\{V_\beta\}_{\beta\in B}$ of $X$
    such that
    \begin{enumerate}[(1)]
        \ii for all $\beta\in B$, there exists $\alpha\in A$ satsifying
	$V_\beta\subseteq U_\alpha$.
	\ii for all $x\in X$, there is a neighborhood $U$ of $x$
	such that $U$ intersects finitely many elements of $\{V_\beta\}_{\beta\in B}$.
    \end{enumerate}
\end{definition}

\begin{prop}
    \label{prop:topman-paracompact}
    Topological manifolds are paracompact.
    More strongly, given a topological manifold $X$,
    an open cover $\{U_\alpha\}_{\alpha\in A}$, and
    any basis $\mathcal B$ for the topology on $X$,
    there exists a countable, locally finite
    refinement of $\{U_\alpha\}_{\alpha\in A}$
    consisting of elements of $\mathcal B$.
\end{prop}

\begin{proof}
    Let $K_1$, $K_2$, $K_3$, \dots be a set of compact subsets of $X$
    such that $K_i\subseteq \operatorname{Int}K_{i+1}$ and
    $\bigcup_{i=1}^\infty K_i = X$.
    Then, define the open set $O_i = \operatorname{Int}K_{i+2}\setminus K_{i-1}$
    and the compact set $C_i = K_{i+1}\setminus\operatorname{Int}K_i$.
    (Such a set exists in general for any second-countable locally compact
    Hausdorff space; see Prop. A.60 in Lee.)

    Now, for each $x\in C_i$, there exists a $U_\alpha$ such that $x\in U_\alpha$
    and a basis element $B_x$ containing $x$ such that $B_x\subseteq U_\alpha\cap O_i$.
    Since $C_i$ is compact, we can find a finite subcover of the $B_x$, say
    $\{B_{i,1}, B_{i,2}, \dots, B_{i,N_i}\}$.
    The set $\{B_{i,j} : j\le N_i\}$ is clearly a refinement of
    $\{U_\alpha\}_{\alpha\in A}$, and we will show local finiteness.

    To do this, given $x\in X$, pick $k$ such that $x\in O_k$.
    If $B_{m,*}\cap O_k\ne\varnothing$, then $C_m\cap O_k\ne\varnothing$,
    which implies that $k-1\le m\le k+2$. There are only finitely many $B_{m,*}$
    satisfying this, so $O_k$ intersects finitely many sets in our refinement,
    proving that the refinement is locally finite and thus $X$ is paracompact.
\end{proof}

\begin{prop}
    \label{prop:strongrefinement}
    Let $X$ be a topological manifold and $\{U_\alpha\}_{\alpha\in A}$
    an open cover that is locally finite.
    Then, there exists a locally finite open cover $\{V_\alpha\}_{\alpha\in A}$
    such that for all $\alpha\in A$, $\ol{V_\alpha}\subseteq U_\alpha$.
\end{prop}

\begin{proof}
    \todo{todo; didn't understand the one shown in class}
\end{proof}

\begin{definition}[$C^\infty$ manifold]
    \label{def:smooth-man}
    Let $X$ be an $n$-dimensional topological manifold.
    An \vocab{atlas} on $X$ is an open cover $\{U_\alpha\}_{\alpha\in A}$ of $X$
    with homeomorphisms
    \begin{equation*}
        h_\alpha\colon U_\alpha\to O_\alpha\stackrel{open}{\subseteq}\RR^n
    \end{equation*}
    such that for all $\alpha,\beta\in A$ with $h_\beta\circ h_\alpha\inv$
    is a diffeomorphism between $h_\alpha(U_\alpha\cap U_\beta)$
    and $h_\beta(U_\alpha\cap U_\beta)$. The pairs $(U_\alpha, h_\alpha)$ are called
    \vocab{coordinate charts} of the underlying topological manifold
    and the composed maps $h_\beta\circ h_\alpha\inv$ are
    called \vocab{transition maps}.
    Then, the pair $(X, \{(U_\alpha,h_\alpha)\}_{\alpha\in A})$,
    abbreviated as $(X, \{U_\alpha\}_{\alpha\in A})$, is a
    \vocab{$C^\infty$ manifold} or \vocab{smooth manifold}
    (which we will just shorten to ``manifold'').
\end{definition}

The transition maps should be thought of as a ``local change of coordinates.''

\begin{definition}[Smooth map to $\RR^N$]
    \label{def:smooth-map-RN}
    Let $(X,\{U_\alpha\}_{\alpha\in A})$ be a manifold
    and $f\colon X\to\RR^N$ a function.
    We say that $f$ is (infinitely) differentiale (or smooth, or $C^\infty$) if
    for all $\alpha\in A$, $f\circ h_\alpha\inv : h_\alpha(U_\alpha)\to\RR^N$
    is $C^\infty$.
\end{definition}

\begin{definition}[Smooth map between manifolds]
    \label{def:smooth-map-man}
    Let $(X,\{U_\alpha\}_{\alpha\in A})$ and
    $(Y,\{V_\beta\}_{\beta\in B})$ be manifolds.
    A map $F\colon X\to Y$ is differentiale (or smooth or $C^\infty$) if
    for all $x\in X$, there exists $U_\alpha\ni x$ and
    $V_\beta\ni f(x)$ and an open set $O$ around
    $h_\alpha(x)$ such that $F(h\inv_\alpha(O))\subseteq V_\beta$,
    and the composition
    \begin{equation*}
	O\stackrel{h\inv_\alpha}{\longrightarrow} U_\alpha
	\stackrel{F}{\longrightarrow} V_\beta
	\stackrel{h_\beta}{\longrightarrow} h_\beta(V_\beta)
    \end{equation*}
    is smooth.
\end{definition}

\begin{definition}[Compatibility of atlases]
    \label{def:atlas-compatibility}
    Two atlases $\mathcal U = \{U_\alpha\}_{\alpha\in A}$ and
    $\mathcal V = \{V_\beta\}_{\beta\in B}$ on $X$
    are \vocab{compatible} if $\mathcal U\cup\mathcal V$ is an atlas
    (so transition maps between elements of $\mathcal U$ and $\mathcal V$
    are smooth).
\end{definition}

Given an open subset $W$ of a manifold $X$ and a diffeomorphism
$\psi\colon W\to O\stackrel{\text{open}}\subseteq \RR^n$,
the atlas of $X$ is compatible with the atlas resulting from
adding $(W,\psi)$ to the atlas. This means that given an atlas
$\mathcal U$ of $X$, we can define
\begin{equation*}
    \mathcal U_{\text{max}} = \mathcal U \cup
    \{ (W,\psi) : \psi\text{ is a diffeomorphism from
	    $W\stackrel{\text{open}}\subseteq X$ to
    an open set in $\RR^n$} \}.
\end{equation*}

It is not hard to check that this is indeed an atlas,
and $\mathcal U_{\text{max}}$ is maximal with respect to inclusion,
meaning it is not properly contained by any other atlases on $X$.

\begin{example}
    Some examples of manifolds:
    \begin{itemize}
        \ii Any open subset $U\subseteq\RR^n$ can be made a manifold
	with the atlas $\{(U,\operatorname{id})\}$.
	(What is the corresponding maximal atlas then?)
	\ii The general linear group $GL_n(\RR)$,
	which is open in $\RR^{n^2}$, can be made a manifold
	as stated above. $GL_n(\RR)$ is also of course a group
	where the multiplication map $(A,B)\mapsto AB$ and
	inverse map $A\mapsto A\inv$ are both $C^\infty$.
	This is an example of a \vocab{Lie group}.
    \end{itemize}
\end{example}

\section{August 28}\label{sec:aug28}
\subsection{Some more examples of manifolds}

The sphere $S^n = \{ x\in\RR^{n+1} : \norm{x} = 1\}$ is
a manifold of dimension of $n$ with charts
\begin{align*}
    \varphi_k^{\pm} : B^n_1(0)&\to S^n\\
    x &\mapsto
    (x_1,\dots,x_{k-1},\sqrt{1-\norm{x}^2},x_k,\dots,x_n)
\end{align*}
where $B^n_r(y)$ indicates an open $n$-dimensional ball
centered at $y$ with radius $r$ and $k=1,\dots,n$.
(It is very simple to verify that the transition maps
are diffeomorphisms; try it!)

The sphere can be acted upon by $\ZZ_2$ via
the antipodal map $x\mapsto -x$ (more precisely,
identity element goes to identity homeomorphism and
the other element goes to the antipodal homeomorphism),
and the orbit space $S^n/\ZZ_2$ is what we call \vocab{real projective space},
denoted $\RR\PP^n$; equivalently, this is
the set of one-dimentional subspaces of $\RR^{n+1}$
(since we can identify each subspace with a pair of antipodal points on $S^n$).
This is a manifold of dimension $n$ with the coordinate charts
\begin{align*}
    \varphi_k\colon\RR^n &\to \RR\PP^n \\
    (x_1,\dots,x_{k-1},x_{k+1},\dots,x_{n+1})&\mapsto
    (x_1,\dots,x_{k-1},1,x_{k+1},\dots,x_{n+1})
\end{align*}
for $k=1,\dots,n+1$. We need to check that
the transition maps are diffeomorphisms in this case.
Note that $\varphi_k(\RR^n)$ is the subset of $\RR\PP^n$
whose $k$th coordinate is nonzero, so if $j\ne k$, then
$\varphi_j(\RR^n)\cap\varphi_k(\RR^n)$ is the subset
whose $j$th and $k$th coordinates are both nonzero,
so the map $\varphi_j\circ\varphi_k\inv\colon\varphi_k(\RR^n)\to\varphi_j(\RR^n)$,
which is simply the division-by-$x_j$ map
\begin{equation*}
    (x_1,\dots,1,\dots,x_k,\dots,x_n)\mapsto
    \left(\frac{x_1}{x_k},\dots,\frac{1}{x_k},\dots,1,\dots,\frac{x_n}{x_k}\right),
\end{equation*}
is differentiable because $x_k\ne 0$, as is its inverse because $x_j\ne 0$.

A generalization of projective space is the \vocab{Grassmannian}
$\Gr(k,n)$ for $0\le k\le n$, which is the space of
all $k$-dimensional subspaces of $\RR^n$.
\todo{i don't get how he computed the dim as $k(n-k)$}

One last example is given by taking a function $F\colon\RR^3\to\RR$
and looking at the pre-image $M = F\inv(0)$ ($0$ is of course arbitrary).
Of course, this need not be a manifold, such as in the case
$F = x^2 + y^2 - z^2$ (no chart exists around $(0,0,0)$);
on the other hand, the simple condition that the derivative $(\nabla F)_x$
at any $x\in M$ is nonzero is enough to make $M$ a manifold!
\todo{prove this}

\subsection{Partitions of unity}

The idea of partitions of unity is to write a smooth function $f\colon X\to\RR$ as
a sum of possibly infinitely many smooth functions such that
the sum locally looks like a finite sum around any $x\in X$,
so we are not restricted to finite sums while also not having to worry
about issues of convergence.

The first step of this is to construct a smooth function $\RR^n\to\RR$
that is supported on a closed ball.
\todo{bump f'ns}

\pagebreak

\section{September 2}\label{sec:sep2}

\todo{partitions of unity stuff}

\todo{def mfld w bd}

The geometric idea of a tangent vector
is that we want to take a $C^\infty$ curve
$\gamma\colon(-\eps,\eps)\to X$,
and then the tangent vector to $\gamma$ at $\gamma(0)$ is
$\frac{\dd\gamma}{\dd t}(0)$.
This, however, depends on the choice of coordinates.
We will instead use a coordinate-free, more abstract path
to define tangent vectors.

\begin{definition}[Germs of a manifold]
    \label{def:germs}
    Let $X$ be a manifold. Consider the collection of pairs
    \begin{equation*}
	C_p = \{ (U,f) : p\in U\overset{\text{open}}\subseteq X,
	f\in C^\infty(U)\}
    \end{equation*}
    and define an equivalence relation $\sim$ on $C_p$ by
    $(U,f)\sim(U',f')$ if there is a neighborhood $W$ of $p$
    contained in $U\cap U'$ such that $f \equiv f'$ on $W$.
    The \vocab{germs} at $p$ are the equivalence classes,
    the collection of which we denote by $C^\infty_p(X)$.
\end{definition}

The germs at $p$ form a unital associative algebra
where the operations of addition, multiplication, and scalar multiplication
are given by
\begin{align*}
    [(U,f)] + [(V,g)] &= [(U\cap V, f+g)]\\
    [(U,f)][(V,g)] &= [(U\cap V, fg)]\\
    \lambda[(U,f)] &= [(U, \lambda f)]\\
\end{align*}
where $[(W,h)]$ is the germ containing $h$.

\begin{exercise}
    Check that these operations are well-defined.
\end{exercise}

The upshot is that we can understand local properties of a function $f$
just by looking at its equivalence class, somehow?

\begin{definition}[Tangent vector]
    \label{def:tangent-vector}
    A \vocab{tangent vector} $v$ at $p\in X$ is a linear map
    $V\colon C^\infty_p(X)\to\RR$ such that
    \begin{equation*}
        v(fg) = v(f)g(p) + f(p)v(g)
    \end{equation*}
\end{definition}
\todo{im too tired for this shit}

\section{September 4}\label{sec:sep4}
\todo{there was some stuff with showing that
partials make basis of tangent space}

\begin{definition}[Tangent bundle]
    \label{def:tangent-bundle}
    The tangent bundle $TX$ is the disjoint union
    \begin{equation*}
	\coprod_{x\in X} T_xX
    \end{equation*}
    equipped with a map $\pi\colon TX\to X$
    given by $\pi(T_xX) = x$.
    A \vocab{section} of $TX$ over an open subset $U\subseteq X$
    is a map $\sigma\colon U\to TX$ such that
    $\pi\circ\sigma = \operatorname{Id}_U$,
    i.e. $\sigma(x)\in T_x(X)$ for all $x\in U$.
\end{definition}

If $\varphi\colon U\to\mathscr O\subseteq\RR^n$
gives local coordinates on $X$, then
\begin{equation*}
    \frac{\partial}{\partial x_j}\colon \mathscr O\to T\mathscr O
\end{equation*}
are local sections of $T\mathscr O$ and w ehave a bijection
$\mathscr O\times\RR^n\to T\mathscr O$ by taking
\begin{equation*}
    (x,c_1,\dots,c_n) \mapsto
    \sum_{j=1}^n c_j\frac\partial{\partial x_j}\bigg\vert_x.
\end{equation*}
The inverse $T\mathscr O\to\mathscr O\times\RR^n$
induces local coordinates on $T\mathscr O$, hence on $TX$.

\begin{prop}
    \label{prop:tangent-bundle-is-mfld}
    If $X$ is a manifold then $TX$ is a $C^\infty$ manifold.
\end{prop}

\begin{proof}
    We need to check that the transition maps between the charts
    described above are diffeomorphisms.\todo{do that.}
\end{proof}

\begin{definition}[Tangent vector field]
    \label{def:tangent-vector-field}
    A \vocab{tangent vector field} $v$ on an open subset $U\subseteq X$
    is a $C^\infty$ section $\sigma\colon U\to TU$.
\end{definition}

\section{September 9}\label{sec:sep9}














\end{document}

